\par\SourceLeaf{127}
Es erhellet hieraus, dass auch die Verbindlichkeit II. sich mit ziemlicher Zuverlässigkeit wird zu Gelde anschlagen. lassen, und dass um diese Rechnungen für I. und II. zu führen, herbeigeschafft werden müssen die Bestimmungen von Geburtsjahr und Tag, für
\begin{quote}
die einzelnen jetzt lebenden Witwen,\par
für deren Kinder unter 20 Jahren, wo solche vorhanden sind, wie bei der Frau Hofr. M., der Frau G. J. R. M. und der Frau Prof. H.\par
für die jetzt verheiratheten Mitglieder,\par
für deren Ehefrauen.
\end{quote}
\noindent\hspace*{\parindent}Ad III. Am bedenklichsten muss aber das Unterfangen erscheinen, den jetzigen Geldwerth der Verbindlichkeit der Kasse gegen die künftigen Theilnehmer \textit{in saecula saeculorum} in Zahlen auszudrucken, versteht sich, nach den jetzigen Statuten, und nach der jetzigen Grösse der Pensionen und Beiträge. Und doch ist es \textit{nothwendig}, dass man in den Stand gesetzt werde, sich hiervon einen wenigstens angenäherten Begriff zu machen, denn es handelt sich ja gerade davon, \textit{dass die Stabilität}, nicht von einer demnächst nach Umständen in ihren Einrichtungen abzuändernden Witwenkasse, sondern \textit{von unsrer Witwenkasse nach ihren jetzigen Einrichtungen begutachtet werden soll.} Man wird hierbei natürlich nicht vergessen, dass die Rechnung von gewissen Elementen abhängig bleibt, die theils schon jetzt nur näherungsweise abzuschätzen sind, theils im Laufe der Zeit sehr bedeutende Abänderungen erleiden können. Von solchen Elementen nenne ich zwei, die Höhe des Zinsfusses und die Anzahl der durchschnittlich jährlich hinzutretenden neuen Mitglieder.

Die Höhe des Zinsfusses steht bei einer Anstalt, die nur zu einem sehr kleinen Theile auf fortgehende Beiträge, der Hauptsache nach auf Capitalrente basirt ist und bleiben soll, offenbar mit der Grösse des erforderlichen Capitals in genauem (verkehrtem) Verhältnisse dergestalt, dass wenn z. B. in einem Zeitpunkte die Schenkung eines Capitals von 70000 Thl. gerade zureichten, eine durchschnittlich immer jährlich gleich viel neue Mitglieder annehmende Gesellschaft zu sustentiren bei einem Zinsfuss von 4 p.c., das Herabsinken des Zinsfusses auf \(3\frac12\) p.c. die Erhöhung des Capitals auf 80000 Thl. erfordern würde. Ich halte diesen Umstand in Beziehung \textit{auf die Schwierigkeit der Begutachtung}, gerade für den unerheblichsten. Die Begutachtung kann mehr nicht thun, als die Grösse des Einflusses in ein klares Licht zu stellen, woraus sich die Folge von selbst ergibt, dass nothwendig schon dafür gesorgt werden muss, dass das Capital durchschnittlich jährlich eine angemessene Erhöhung erhalte, um dem im Laufe der Zeit allmählig aber unfehlbar eintretenden Sinken des Zinsfusses zu begegnen.

Ebenso einleuchtend ist es, dass die Grösse des erforderlichen Capitals genau der Anzahl der durchschnittlich jährlich beitretenden neuen Mitglieder (\textit{ceteris paribus}) proportional sein wird. Wir können zunächst nur unsre eignen Erfahrungen zum Grunde legen, die seit 100 Jahren vorliegen, und wo natürlich die neuern und neuesten Zeiten unser Urtheil vorzugsweise leiten müssen. Se. Magnificenz bemerkt mit Recht, dass die aus den gesteigerten wissenschaftlichen Bedürfnissen und Anforderungen hervorgegangene Vergrösserung der Zahl der Professoren einen wesentlichen Einfluss auf das Bestehen solcher Professorenwitwenkassen haben muss, die hauptsächlich auf Capital fundirt sind. Es ist also sehr wohl möglich, dass die Göttingischen Ergebnisse z. B. seit den letzten 30 oder 40 Jahren keinen ganz richtigen Maassstab für die Zukunft, zumal für die Zukunft späterer Jahrhunderte bilden können; aber diese Ungewissheit liegt in der Natur der Veränderlichkeit aller menschlichen Dinge, die Folgen davon treten \textit{allmählig} hervor, und man begegnet ihnen nur durch eine niemals einschlummernde Vigilanz. In unserm Falle also macht man seine Rechnung für das zur nachhaltigen Erfüllung der Verbindlichkeit III. erforderliche Capital nach unsern besten jetzigen Kenntnissen, vergisst nicht, dass eben wegen jener Ungewissheit ein etwas grösseres Capital vorhanden sein müsse, \textit{wiederholt} die Rechnung fortwährend in bestimmten nicht gar zu grossen
